% ===== wrap-up =====
\begin{frame}{이번 주차 정리}
  \begin{enumerate}
    \item \kb{9.1 --- 함수근사의 진짜 이유는 일반화} ---
      Atari 상태 공간 $10^{16{,}993}$은 표로 원리적으로 불가.
      ``4,610개 파라미터가 10억 칸을 대체''의 핵심은
      \textbf{보지 못한 상태에 대한 보간}.
      손실의 $\max_{a'} Q(s',a';\theta)$ 항이 ``목표가 움직인다''의 근원.
    \item \kb{9.2 --- 타겟 네트워크} ---
      목표 계산 전용 $\theta^{-}$를 $C$스텝마다 하드 동기화해
      ``움직이는 과녁''을 고정. 실험: 타겟 없이는
      \textbf{+27\% 과대추정}로 ``잘못된 값에 안정''.
      낮은 TD 손실 $\ne$ 올바른 Q.
    \item \kb{9.3 --- 경험 재현} ---
      버퍼에서 무작위 미니배치로 \textbf{상관 제거 $+$ 데이터 재사용}.
      off-policy가 낡은 경험 재사용을 허용.
      CartPole 21$\to$104(시드 0), 세 시드 모두 만점 500 도달,
      \textbf{여러 시드의 평균 곡선}으로 판단.
  \end{enumerate}
  \vskip0.8em
  \begin{block}{다음 주 --- Week 10. 정책기반 강화학습}
    DQN은 매 스텝 $\max_{a'} Q(s',a')$를 계산하는 구조.
    \textbf{행동이 연속량이면} 이 최댓값을 취하는 것 자체가 불가능
    $\Rightarrow$ Q-함수를 우회하고 \textbf{정책을 직접 학습}하는
    policy-based RL(REINFORCE)로 건너든다.
  \end{block}
\end{frame}
